% Chapter 3: Context — Handbook 8.2.4
% Literature review organised BY TOPIC, not by paper.
% ~1000 words for BSc. Cite everything. Show relevance to your project.

\section{Context}
\label{sec:context}

\subsection{Secure Multi-Party Computation}
\label{sec:ctx-mpc}

MPC enables a set of parties to jointly compute a function over their private
inputs without revealing those inputs to one
another~\cite{evans2018_pragmatic_mpc}.  The following definitions use the
formalisations in Evans et al.~(\citeyear{evans2018_pragmatic_mpc}) and
Zhao et al.~(\citeyear{zhao2019_secure_mpc}).  Two security models are central
to this project.

\subsubsection{Shamir Secret Sharing (Honest Majority)}
\label{sec:ctx-shamir}

Shamir's Secret Sharing~\cite{shamir1979_secret_sharing} hides a secret $s$
inside a random polynomial of degree $t$ whose constant term is $s$.  Each
party receives the polynomial evaluated at a distinct point.  Any $t{+}1$
parties can reconstruct $s$ via interpolation, but $t$ or fewer learn nothing
about it.

When used for MPC, the threshold is set to $t < N/2$ (honest majority).
Addition of shared values is free (each party adds its shares locally),
while multiplication requires one round of communication per gate, giving
$O(N)$ messages per multiplication.  This makes Shamir efficient at large
party counts.  However, if the honest-majority assumption is violated
($t \geq N/2$ parties are corrupt), the corrupt parties can reconstruct
every shared secret, and all privacy guarantees are lost.

\subsubsection{MASCOT (Dishonest Majority)}
\label{sec:ctx-mascot}

The MASCOT protocol~\cite{keller2016_mascot} is built on the SPDZ
framework~\cite{damgaard2011_spdz} and provides malicious security tolerating
up to $N{-}1$ corrupt parties.  The MPC program is expressed as an arithmetic
circuit: each operation (addition, multiplication, comparison) is a gate, and
the parties jointly evaluate the circuit by exchanging messages at each gate.
MASCOT operates in two phases.

\textbf{Offline phase (preprocessing).}  Before any inputs are known, the
parties generate a pool of authenticated Beaver triples using oblivious
transfer (OT).  For each pair of parties $(i, j)$, a series of OT instances
produces correlated random values that, when combined across all pairs, yield
a globally shared triple $(a, b, c)$ with $c = a \cdot b$.

Every shared value is authenticated with a MAC under a global key $\alpha$
that is itself secret-shared among all parties~\cite{damgaard2011_spdz}.
No single party knows $\alpha$, so no party can forge a valid share.

The pairwise nature of OT means the offline phase requires $O(N^2)$
communication: every pair of parties must exchange OT messages.  This is the
dominant cost and the reason MASCOT is significantly more expensive than
Shamir, particularly at higher party counts.

\textbf{Online phase (computation).}  Addition of shared values is free
(each party adds locally).  Multiplication consumes one pre-computed Beaver
triple per gate: the parties open masked values that reveal nothing about the
inputs, then combine the result with the triple to obtain a sharing of the
product.

The key difference from Shamir is what happens when a value is \emph{opened}
(revealed to all parties).  Before accepting an opened value $v$, all parties
perform a MAC check: each party $i$ computes $m_i - v \cdot \alpha_i$ and
the parties jointly verify that these sum to zero.  If any corrupt party
modified their share, the MAC check fails with overwhelming probability and
the protocol aborts.  This is what provides malicious security: corrupt
parties can deviate arbitrarily from the protocol, but any cheating is
detected before it affects the output.

A third protocol, replicated secret sharing, splits each secret into three
additive shares and gives each of the three parties two of them.  It is the
simplest multi-party scheme but is restricted to exactly three parties with an
honest majority.  It was used during early development to validate the MPC
integration pipeline before moving to MASCOT and Shamir for evaluation.

The MP-SPDZ framework~\cite{mp-spdz} provides academic production-grade
implementations of all three protocols, allowing the same high-level auction
program to be compiled and executed under different protocol backends.  This
was central to the evaluation in Chapter~\ref{sec:results}, where the same
program was benchmarked under MASCOT and Shamir with identical inputs.  Other
MP-SPDZ backends (semi-honest two-party, BMR garbled circuits, etc.) were not
tested because they either do not support $N > 2$ parties or solve a different
threat model than the one defined in Section~\ref{sec:intro-security-policy}.

The Danish sugar-beet auction~\cite{bogetoft2009_mpc_live} remains the most
prominent real-world MPC deployment for auctions, using a three-party
honest-majority protocol.  That system relied on a small number of known,
trusted intermediaries and a direct TCP network.  This project extended the
concept to an untrusted, peer-to-peer setting where the honest-majority
assumption may not hold, motivating the comparison between MASCOT and Shamir.

Zhao et al.~\cite{zhao2019_secure_mpc} survey MPC theory and applications.
Feng and Yang~\cite{feng2022_concretely_efficient_mpc} provide concrete
efficiency benchmarks, noting that the choice between honest and dishonest
majority remains the dominant factor in practical performance.  Li et
al.~\cite{li2023_secure_comp_tee_survey} compare MPC with trusted execution
environments (TEEs), concluding that MPC provides stronger trust assumptions
at the cost of higher communication overhead.

\subsection{Veilid: Privacy-Preserving Peer-to-Peer Networking}
\label{sec:ctx-veilid}

Veilid is described by its developers as
\begin{quote}
``an open-source, peer-to-peer application framework designed to build
privacy for all''~\cite{veilid_website}.
\end{quote}
Each application embeds a Veilid node into a global overlay where peers are
equal (no privileged relays), connections are end-to-end encrypted, and
private routing obscures network locations.

Three capabilities made Veilid suitable for this project:
\begin{enumerate}[noitemsep]
  \item \textbf{Addressable encrypted endpoints} keyed by public keys rather
    than IP addresses, eliminating direct IP exposure.
  \item \textbf{A secure DHT} for publishing listings, bids, and discovering
    auction participants.
  \item \textbf{Onion-routed private routes} for point-to-point messaging
    that limit metadata exposure during bid submission and MPC tunnel traffic.
\end{enumerate}

\subsubsection{What Veilid Protects (and What It Does Not)}
\label{sec:ctx-veilid-security}

Veilid's private routes provide network-layer anonymity: a sender constructs
a multi-hop route through relay nodes, and the receiver sees only the last
hop's identity.  This protects participants' IP addresses from each other and
from passive network observers.  However, several limitations were relevant to
the security policy:

\begin{itemize}[noitemsep]
  \item \textbf{Anonymity set}: Privacy depends on the number of available
    relay nodes.  On a small devnet (20--80 nodes), the anonymity set is
    correspondingly small.  On the public Veilid network this would be larger
    but remains unquantified.
  \item \textbf{DHT metadata}: Writing a bid record to the DHT is observable
    by DHT nodes.  The \emph{number} and \emph{timing} of bids are visible
    even though \emph{values} are encrypted.
  \item \textbf{Route liveness}: Private routes are application-managed and
    can die silently.  There is no notification when a remote route becomes
    invalid; sending to a dead route fails only after a timeout.  This had
    significant implications for MPC integration
    (Section~\ref{sec:results-mpc}).
  \item \textbf{IP leakage without private routes}: If a node communicates
    using unsafe routing (direct node-to-node), its IP address is exposed on
    the wire.  All benchmarked messaging in this system used safe routing
    exclusively.
\end{itemize}

Table~\ref{tab:overlay-comparison} compares Veilid with three alternative
overlay networks that offer similar privacy properties.

\begin{table}[H]
\centering
\small
\renewcommand{\arraystretch}{1.3}
\makebox[\textwidth][c]{%
\begin{tabular}{|l|l|l|l|l|}
\hline
\textbf{Property} & \textbf{Tor} & \textbf{I2P} & \textbf{IPFS} & \textbf{Veilid} \\
\hline
Primary model     & Client--server (exit nodes) & Internal services & Content-addressed storage & P2P application framework \\
\hline
Anonymity routing & Onion circuits (3 hops)     & Garlic routing     & None (cleartext DHT)      & Onion-routed private routes \\
\hline
Built-in DHT      & Hidden-service descriptors  & LeaseSet/NetDB     & Content routing            & General-purpose key--value \\
\hline
Embeddable library & No (SOCKS proxy)           & Partial (I2CP)     & Yes (libp2p)               & Yes (veilid-core) \\
\hline
P2P messaging     & Via hidden services         & Via streaming lib   & No                         & \texttt{app\_call} / \texttt{app\_message} \\
\hline
\end{tabular}}%
\caption{Overlay network comparison.}
\label{tab:overlay-comparison}
\end{table}

Tor~\cite{dingledine2004_tor} provides strong anonymity through onion-routed
circuits but is designed for the client--server model: clients reach clearnet
servers via exit nodes, or hidden services via rendezvous points.  Building a
peer-to-peer auction on Tor would require running each participant as a hidden
service, managing circuit lifecycle externally, and implementing DHT storage
separately.  Tor's hidden services also impose a minimum of six hops per
connection (three per side), adding latency that compounds across MPC rounds.

I2P~\cite{zantout2011_i2p} is closer to a P2P model, using garlic routing to
bundle messages for internal (``eepsite'') services.  Its streaming library
supports bidirectional TCP-like connections between I2P destinations.  However,
I2P lacks a general-purpose DHT for application data (its NetDB stores only
router and destination metadata), and its ecosystem is smaller with fewer
maintained client libraries.

IPFS~\cite{benet2014_ipfs} provides content-addressed storage over a
Kademlia-based DHT and is well suited for data distribution.  However, it
provides no anonymity: DHT operations and data transfers expose participants'
IP addresses.  Adding a privacy layer would require tunnelling IPFS through Tor
or I2P, combining the limitations of both.

Veilid was selected because it provides all three capabilities required by the
application in a single embeddable library: a general-purpose DHT for
storing listings and bids, onion-routed private routes for anonymous messaging,
and a peer-to-peer architecture where all nodes are equal with no privileged
relays or exit nodes.  This avoided the need to compose multiple
frameworks (e.g.\ IPFS for storage plus Tor for anonymity) and the integration
complexity that would entail.

\subsubsection{Bootstrap and Network Formation}
\label{sec:ctx-veilid-bootstrap}

A Veilid network requires at least one \emph{bootstrap node}: a
well-known server whose address is configured into every joining node so
that they can discover initial peers and populate their routing
tables~\cite{veilid_bootstrap_setup}.  On the public Veilid network,
bootstrap nodes are discovered via DNS TXT records.  For development and
testing, a private devnet must provide its own bootstrap infrastructure.

Crucially, Veilid expects each node to present a unique public IP address.
Its routing-domain logic classifies incoming connections by IP range and
drops traffic from addresses it considers local or
duplicate~\cite{veilid_developer_book}. Veilid also implements its own STUN for nodes
behind some level of NAT, but nodes behind NAT are not deemed ``optimal'', and recieve
network connection penalties accordingly (cannot act as bootstrap/less priority for relay).
This means that running multiple nodes on a single machine as required for a local devnet fails
unless each node appears to be on the public internet and not behind a NAT layer.  A community-contributed development
network~\cite{veilid-devnet-mr444} uses Docker containers with
per-container IP addresses.  This project's playground devnet took a lighter
approach: \texttt{LD\_PRELOAD} IP spoofing via a shared library
(\texttt{libipspoof.so}) that intercepts socket calls and maps each process to
a unique fake global IP (e.g.\ \texttt{1.2.3.1} through \texttt{1.2.3.255}),
supporting up to 255 nodes on a single machine without containerisation.
The implications of this constraint for devnet design are discussed further in
Section~\ref{sec:method-approach}.

\subsection{Stigmerge and the Route-in-DHT Pattern}
\label{sec:ctx-stigmerge}

Stigmerge~\cite{stigmerge2024} is a decentralised file-sharing application
built on Veilid that demonstrated the coordination patterns this project
adapted.  The name refers to \emph{stigmergy},indirect coordination through
a shared environment,where peers communicate by reading and writing shared
DHT records rather than maintaining direct connections.

Stigmerge established two patterns adopted in this project:
\begin{enumerate}[noitemsep]
  \item \textbf{Route-in-DHT}: A node creates a private route, serialises the
    route blob, and writes it to a DHT subkey.  Remote peers read the subkey,
    import the blob, and send messages via \texttt{app\_call} to the imported
    route.  This decouples route exchange from direct messaging.
  \item \textbf{\texttt{app\_call} for confirmed delivery}: All point-to-point
    data transfer uses Veilid's \texttt{app\_call} (request--response/SYN \& ACK) rather
    than \texttt{app\_message} (fire-and-forget), ensuring that failed
    deliveries are detected immediately rather than silently dropped.
\end{enumerate}

These patterns were critical for MPC tunnelling, where every message had to be
confirmed and route failures had to be detected promptly.

\subsection{Auction Privacy in Context}
\label{sec:ctx-auction-privacy}

Sako~\cite{sako2000_hide_losers} was the first to propose a cryptographic
auction protocol in which losing bids are never revealed.  The scheme uses
homomorphic encryption and zero-knowledge proofs so that the auctioneer can
publicly verify that the announced winner holds the highest bid, without
disclosing the values submitted by any other participant.  Sako proved the
protocol's correctness and privacy properties formally but did not produce an
implementation, leaving practical feasibility as an open question.
Blockchain-based auctions remove
the trusted auctioneer via smart contracts but expose all bids on-chain and
invite front-running~\cite{daian2019_flashboys2}.  The Danish sugar-beet
auction~\cite{bogetoft2009_mpc_live} demonstrated MPC in production but relied
on three trusted intermediaries operating over direct TCP.

This project combined the privacy guarantees of MPC (non-winning bids never
revealed) with the transport-layer anonymity of Veilid (participant IPs never
exposed), in a fully decentralised setting with no trusted auctioneer.

The seller occupies a semi-trusted role as MPC party~0.  This is justified
on two grounds.  First, the seller has a direct economic incentive to complete
the auction honestly: withholding the decryption key or manipulating results
means the sale does not proceed, costing the seller revenue.  Second, the
seller is the inherent key holder: since they encrypted the listing content,
they necessarily possess the decryption key, and any scheme that removes
the seller from key delivery requires threshold cryptography
(Section~\ref{sec:conc-future}).  The seller's power is strictly limited
ephemeral route to the winner, not all bids or bidder identities.

Chapter~\ref{sec:results} evaluates the cost of this combination in terms of
performance overhead, reliability, and residual information leakage.


%% TODO: Architecture diagrams for Ch5 (sequence, module dependency)